\section{Transformer 的应用与影响}

\subsection{Attention Is All You Need}

2017 年，Vaswani 等人发表了划时代的论文 \textit{Attention Is All You Need}，提出了 Transformer 架构。该架构\textbf{完全摒弃了循环机制}，仅基于 Self-Attention 和前馈层构建序列模型。

\begin{figure}[H]
\centering
\includegraphics[width=0.6\textwidth]{slides-images/slide-65.jpg}
\caption{``Attention Is All You Need''——Transformer 时代的开启\protect\footnotemark}
\end{figure}
\footnotetext{Slides 第65页。}

\begin{knowledgebox}{GPT 中的 T 代表 Transformer}
如果你使用过 ChatGPT，那么名字中的 ``T'' 就代表 Transformer。GPT 全称是 Generative Pre-trained Transformer——一种基于 Transformer 架构的生成式预训练模型。Transformer 是当今几乎所有大语言模型的底层架构。
\end{knowledgebox}

\subsection{Transformer vs. RNN：优势对比}

\begin{table}[H]
\centering
\caption{RNN 与 Transformer 的对比}
\begin{tabular}{lcc}
\toprule
\textbf{特性} & \textbf{RNN/LSTM} & \textbf{Transformer} \\
\midrule
处理方式 & 逐时间步（Sequential） & 全序列并行（Parallel） \\
长期依赖 & 困难（梯度消失） & 直接建模（Self-Attention） \\
训练效率 & 低（无法并行化） & 高（可充分利用 GPU） \\
位置信息 & 隐式（通过时间步顺序） & 显式（Positional Encoding） \\
可扩展性 & 有限 & 强（支持超大规模） \\
\bottomrule
\end{tabular}
\end{table}

\subsection{Self-Attention 的广泛应用}

Self-Attention 和 Transformer 的影响远超自然语言处理：

\begin{figure}[H]
\centering
\includegraphics[width=0.85\textwidth]{slides-images/slide-80.jpg}
\caption{Self-Attention 在不同领域的应用\protect\footnotemark}
\end{figure}
\footnotetext{Slides 第80页。}

\begin{itemize}
    \item \textbf{自然语言处理}：BERT（Devlin et al., NAACL 2019）、GPT（Brown et al., NeurIPS 2020）等大语言模型
    \item \textbf{生物序列}：AlphaFold（Jumper et al., Nature 2021）用 Attention 预测蛋白质三维结构；ESMFold（Lin et al., Science 2023）
    \item \textbf{计算机视觉}：Vision Transformer / ViT（Dosovitskiy et al., ICLR 2020）将图像切分为 Patch 序列，用 Transformer 处理，性能媲美甚至超越卷积神经网络
    \item \textbf{多模态}：跨模态的 Attention 机制实现了图文生成（如 DALL-E）、视频理解等任务
\end{itemize}

\begin{importantbox}{Transformer 的通用性}
Transformer 的成功关键在于其\textbf{与领域无关}（Domain-agnostic）的架构设计。无论输入是文字、图像、蛋白质序列还是音频，只要能将其表示为一系列向量（Token），就可以用 Transformer 进行处理。这种通用性使得 Transformer 成为了``基础模型''（Foundation Model）时代的标准架构。
\end{importantbox}

\subsection{本章小结}

Transformer 彻底改变了序列建模的格局。它通过 Self-Attention 机制消除了对循环结构的依赖，实现了高效的并行计算和直接的长距离依赖建模。从自然语言处理到蛋白质结构预测，从计算机视觉到多模态生成，Transformer 的影响无处不在。

\newpage
